% block: matrix-exponential-heading page:210 type:section
\subsection{{Matrix Exponential Function}}\label{sec:matrix-exponential-function}

% block: matrix-exponential-series-definition page:210 type:ordinary-paragraph
{As seen in Definition~\ref{def:jordan-matrix-function}, the matrix exponential function is defined by the Jordan canonical form. On the other hand, the following equivalent definition is useful for computing the matrix function:}
\begin{equation}\label{eq:matrix-exponential-series}e^A:=I+A+\frac{1}{2!}A^2+\frac{1}{3!}A^3+\cdots.\end{equation}
{The series converges for any square matrix $A\in\mathbb{C}^{n\times n}$. In this section, numerical algorithms for matrix exponential functions are described.}

% block: matrix-exponential-properties-transition page:210 type:ordinary-paragraph
{Fundamental properties of matrix functions are summarized in Theorem~\ref{thm:matrix-exponential-properties}. These properties are easily proved from the definition in \eqref{eq:matrix-exponential-series}.}

% block: matrix-exponential-properties-theorem page:210 type:theorem
\Needspace{14\baselineskip}
\begin{Theorem}{{Properties of matrix exponential functions}}\label{thm:matrix-exponential-properties}
{For $A,B\in\mathbb{C}^{n\times n}$:}
\begin{enumerate}[label=(\arabic*)]
\item\label{part:matrix-exp-zero} {$e^O=I$, where $O$ is the zero matrix;}
\item\label{part:matrix-exp-commuting-sum} {$e^Ae^B=e^{A+B}$ if $AB=BA$;}
\item\label{part:matrix-exp-inverse} $e^Ae^{-A}=I$, $(e^A)^{-1}=e^{-A}$;
\item\label{part:matrix-exp-similarity} {$e^{XAX^{-1}}=Xe^AX^{-1}$, where $X$ is a nonsingular matrix;}
\item\label{part:matrix-exp-adjoint} $e^{A^{\mathrm H}}=(e^A)^{\mathrm H}$.
\end{enumerate}
\end{Theorem}

% block: matrix-exponential-hermitian-unitary page:210 type:ordinary-paragraph
{From \ref{part:matrix-exp-adjoint} in Theorem~\ref{thm:matrix-exponential-properties}, matrix exponential function $e^A$ is Hermitian if $A$ is Hermitian. From \ref{part:matrix-exp-inverse} and \ref{part:matrix-exp-adjoint} in Theorem~\ref{thm:matrix-exponential-properties}, matrix exponential function $e^A$ is unitary, i.e., $(e^A)^{\mathrm H}e^A=e^{-A}e^A=I$ if $A$ is skew-Hermitian $(A^{\mathrm H}=-A)$.}

% block: matrix-exponential-scalar-multiples page:211 type:ordinary-paragraph
{From \ref{part:matrix-exp-commuting-sum} in Theorem~\ref{thm:matrix-exponential-properties} with $A=aC$ and $B=bC$ for scalar values $a$ and $b$, we have $e^{aC}e^{bC}=e^{(a+b)C}$. Here, if $a=1$ and $b=-1$, then we have $e^Ce^{-C}=e^{0C}=e^O=I$ that corresponds to \ref{part:matrix-exp-inverse} in Theorem~\ref{thm:matrix-exponential-properties}.}

% block: matrix-exponential-algorithms-heading page:211 type:subsection
\subsubsection{{Numerical Algorithms for Matrix Exponential Functions}}\label{subsec:matrix-exponential-algorithms}

% block: matrix-exponential-direct-truncation page:211 type:ordinary-paragraph
{One of the simplest ways to approximately compute matrix exponential functions is to truncate the series in \eqref{eq:matrix-exponential-series}, i.e., $I+A+\frac{1}{2!}A^2+\frac{1}{3!}A^3+\cdots+\frac{1}{m!}A^m$. If all the eigenvalues of matrix $A$ are close to zero, then the truncation error may be small for small $n$. The above approach is, however, inefficient for a matrix $A$ with no such distribution of the eigenvalues of $A$.}

% block: matrix-exponential-scaled-truncation page:211 type:ordinary-paragraph
{A better and simple way to compute matrix exponential functions is to consider $e^{A/s}$ instead of $e^A$. We see that $(e^{A/s})^s=e^A$ for $s\ge1$, and the all eigenvalues of $A/s$ are $s$ times smaller than the eigenvalues of $A$. Then truncating the following series:}
\begin{equation}\label{eq:matrix-exp-scaled-series}e^{A/s}=I+\frac{1}{1!s}A+\frac{1}{2!s^2}A^2+\frac{1}{3!s^3}A^3+\cdots\end{equation}
{yields}
\begin{equation}\label{eq:matrix-exp-scaled-truncation}F_{r,s}=I+\frac{1}{1!s}A+\frac{1}{2!s^2}A^2+\frac{1}{3!s^3}A^3+\cdots+\frac{1}{r!s^r}A^r,\end{equation}
{which is an approximation to $(F_{r,s})^s\approx e^A$. Then, the following property holds:}\footnote{{In \BookCite{179}, Theorem~\ref{thm:matrix-exp-truncation-error} is stated on Banach algebra, which is a more general result.}}

% block: matrix-exp-truncation-error-theorem page:211 type:theorem
\Needspace{14\baselineskip}
\begin{Theorem}{\BookCite{179}}\label{thm:matrix-exp-truncation-error}
{Let $A\in\mathbb{C}^{n\times n}$ and $F_{r,s}$ be the truncation given in \eqref{eq:matrix-exp-scaled-truncation}. Then}
\begin{equation}\label{eq:matrix-exp-truncation-error}\|e^A-(F_{r,s})^s\|\le\frac{\|A\|^{r+1}e^{\|A\|}}{s^r(r+1)!},\end{equation}
{where $\|\cdot\|$ is any matrix norm that is submultiplicative.}\footnotemark
\end{Theorem}
\footnotetext{{See (\ref{cond:matrix-norm-product}) in Section~\ref{subsec:matrix-norm} for the term ``submultiplicative''.}}

% block: matrix-exp-truncation-parameter-example page:211 type:ordinary-paragraph
{From the right-hand side of Theorem~\ref{thm:matrix-exp-truncation-error}, the truncation error can be estimated by computing $\|A\|$ or its upper bound. For example, we consider the case $s=2^m$. Assume that $\|A\|=10$. Then if we use $r=8$, $m=9$, $s=2^m$, the right-hand side of \eqref{eq:matrix-exp-truncation-error} is about $1.29\times10^{-14}$, which means that it is possible to determine $s$ and $r$ such that the truncation error is less than a given tolerance, if $\|A\|$ is estimated. When $\|A\|$ is a Frobenius norm, then it is easy to compute the norm by the definition of a Frobenius norm. When $\|A\|$ is a matrix 2-norm, the maximum singular value of $A$ is needed, which can also be easily estimated by the Golub--Kahan bidiagonalization process (Algorithm~\ref{alg:golub-kahan}).}

% block: matrix-exp-pade-approximant page:212 type:ordinary-paragraph
{A more cost-efficient method is as follows: instead of computing \eqref{eq:matrix-exp-scaled-truncation}, we use the Padé approximant to $e^{A/s}$ and compute $s$th power of the approximant. The Padé approximant $R_{p,q}(A)$ to matrix exponential function $e^A$ is given below:}
\begin{equation}\label{eq:matrix-exp-pade}R_{p,q}(A)=[D_{p,q}(A)]^{-1}N_{p,q}(A),\end{equation}
{where}
\[\begin{aligned}N_{p,q}(A)&=\sum_{j=0}^p n_jA^j,\quad n_j=\frac{(p+q-j)!}{(p+q)!}\times\frac{p!}{j!(p-j)!},\\D_{p,q}(A)&=\sum_{j=0}^q d_j(-A)^j,\quad d_j=\frac{(p+q-j)!}{(p+q)!}\times\frac{q!}{j!(q-j)!}.\end{aligned}\]

% block: matrix-exp-diagonal-pade-parameters page:212 type:ordinary-paragraph
{When approximating the matrix exponential, it is recommended to use the diagonal Padé approximant, i.e., $p=q$. From $e^A=[e^{A/2^m}]^{2^m}\approx[R_{q,q}(A/2^m)]^{2^m}$, we need to choose two parameters $q$ and $m$. If $\|A\|/2^m\le1/2$, then}
\[[R_{q,q}(A/2^m)]^{2^m}=e^{A+E},\quad\frac{\|E\|}{\|A\|}\le8\left[\frac{\|A\|}{2^m}\right]^{2q}\frac{(q!)^2}{(2q)!(2q+1)!}.\]
{For the details, see [\hyperref[bib:source-133]{\ref*{bib:source-133}}, p.12]. The inequality leads to optimum parameters $q$ and $m$ such that $q+m$ is minimized under the condition that the truncation error holds $\|E\|/\|A\|\le\epsilon$ for a given tolerance $\epsilon$. The optimum parameters for a tolerance $\epsilon=10^{-15}$ with respect to $\|A\|$ are given in Table~\ref{tab:matrix-exp-pade-parameters}. For related studies, see \BookCite{8}, \BookCite{98}.}

% block: matrix-exp-schur-reduction page:212 type:ordinary-paragraph
{If Schur decomposition $A=QTQ^{\mathrm H}$ is computed, then it follows from \eqref{eq:matrix-exponential-series} that $e^A=e^{QTQ^{\mathrm H}}=Qe^TQ^{\mathrm H}$. Thus all we have to do is to compute $e^T$ using (diagonal) Padé approximant.}

% block: matrix-exp-hermitian-eigen-decomposition page:212 type:ordinary-paragraph
{In particular, when $A$ is Hermitian, $T$ becomes a diagonal matrix. Let $\lambda_1,\lambda_2,\ldots,\lambda_n$ be the diagonal elements of $T$. Then the matrix exponential function can be computed by}
\begin{equation}\label{eq:matrix-exp-hermitian-diagonal}e^A=Q\begin{bmatrix}e^{\lambda_1}&&\\&\ddots&\\&&e^{\lambda_n}\end{bmatrix}Q^{\mathrm H}.\end{equation}
{Further, if $A$ is diagonalizable, i.e., $A=XDX^{-1}$ for a nonsingular matrix $X$ and a diagonal matrix $D$, then we have $e^A=Xe^DX^{-1}$, where $e^D$ is the same form as in \eqref{eq:matrix-exp-hermitian-diagonal}. On the other hand, if the condition number $X$ is high, this approach can be numerically unstable.}

% block: matrix-exp-pade-parameter-table page:212 type:table
\begin{table}[H]\centering
\caption{{Optimum parameters $q$ and $m$ in $R_{q,q}(A/2^m)$ satisfying the error tolerance $\epsilon=10^{-15}$.}}\label{tab:matrix-exp-pade-parameters}
\begin{tabular}{c|cccccccccc}\hline
$\|A\|$&$10^{-3}$&$10^{-2}$&$10^{-1}$&$1$&$10$&$10^2$&$10^3$&$10^4$&$10^5$&$10^6$\\\hline
$q$&3&3&4&6&6&6&6&6&6&6\\
$m$&0&0&0&1&5&8&11&15&18&21\\\hline
\end{tabular}\end{table}

% block: matrix-exponential-vector-heading page:213 type:subsection
\subsubsection{{Multiplication of a Matrix Exponential Function and a Vector}}\label{subsec:matrix-exponential-vector}

% block: matrix-exponential-ode-solution page:213 type:ordinary-paragraph
{We consider the following linear ordinary differential equation:}
\begin{equation}\label{eq:matrix-exp-linear-ode}\frac{\mathrm{d}\boldsymbol{x}(t)}{\mathrm{d}t}=A\boldsymbol{x}(t)+\boldsymbol{f}(t),\quad\boldsymbol{x}(0)=\boldsymbol{x}_0,\end{equation}
{where $A\in\mathbb{C}^{n\times n}$, $\boldsymbol{x},\boldsymbol{f}\in\mathbb{C}^n$, $t\in\mathbb{R}$. The solution can be given by using matrix exponential functions as follows:}
\begin{equation}\label{eq:matrix-exp-variation-of-constants}\boldsymbol{x}(t)=e^{tA}\boldsymbol{x}_0+\int_0^t e^{(t-\tau)A}\boldsymbol{f}(\tau)\,\mathrm{d}\tau.\end{equation}

% block: matrix-exponential-vector-product-motivation page:213 type:ordinary-paragraph
{We see that in \eqref{eq:matrix-exp-variation-of-constants} there are multiplications of matrix exponential functions and vectors. In particular, if $\boldsymbol{f}(t)=\boldsymbol{0}$, then the solution is written as $\boldsymbol{x}(t)=e^{tA}\boldsymbol{x}_0$. This explicit form of the solution means that we know the solution at an arbitrary time $t$ if the multiplications are obtained.}

% block: matrix-exponential-vector-cost page:213 type:ordinary-paragraph
{If matrix exponential functions have already been computed, then all we have to do is to compute the matrix--vector multiplication whose computational cost is of order $n^2$, where $n$ is a matrix of size $n\times n$. On the other hand, computing the cost of a matrix exponential function may be of order $n^3$ even if matrix $A$ is sparse. Thus computing matrix exponential functions is inefficient, when it comes to computing the solution of \eqref{eq:matrix-exp-linear-ode}. Note that the solution only requires the result of the multiplication of a matrix exponential functions and a vector. In such a case, an approximation to the multiplication using Krylov subspace is useful, which is described below.}

% block: matrix-exponential-krylov-least-squares page:213 type:ordinary-paragraph
{Since matrix exponential function $e^{tA}$ is defined by the series \eqref{eq:matrix-exponential-series}, it is natural to use matrix polynomials $p_m(A)$ of degree $m$ to approximate $e^{tA}$. The multiplication of a matrix polynomial $p_m(A)$ and a vector belongs to Krylov subspace $\mathcal{K}_m(A,\boldsymbol{x}_0)$. Thus the optimal approximation is given by solving the following least-squares problem:}
\begin{equation}\label{eq:matrix-exp-krylov-least-squares}\min_{\boldsymbol{x}_m\in\mathcal{K}_m(A,\boldsymbol{x}_0)}\|e^{tA}\boldsymbol{x}_0-\boldsymbol{x}_m\|.\end{equation}

% block: matrix-exponential-orthogonal-projection page:213 type:ordinary-paragraph
{Let $V_m$ be a matrix whose columns are orthonormalized basis vectors of $\mathcal{K}_m(A,\boldsymbol{x}_0)$. Then, any vector $\boldsymbol{x}\in\mathcal{K}_m(A,\boldsymbol{x}_0)$ can be written as $\boldsymbol{x}=V_m\boldsymbol{y}$, where $\boldsymbol{y}\in\mathbb{C}^m$. Thus \eqref{eq:matrix-exp-krylov-least-squares} is equivalent to}
\[\min_{\boldsymbol{y}\in\mathbb{C}^m}\|e^{tA}\boldsymbol{x}_0-V_m\boldsymbol{y}\|.\]
{The solution is $\boldsymbol{y}=V_m^{\mathrm H}e^{tA}\boldsymbol{x}_0$, and thus the solution of \eqref{eq:matrix-exp-krylov-least-squares} is given by}
\begin{equation}\label{eq:matrix-exp-krylov-projection}(\boldsymbol{x}(t)\approx)\ \boldsymbol{x}_m=V_mV_m^{\mathrm H}e^{tA}\boldsymbol{x}_0.\end{equation}

% block: matrix-exponential-arnoldi-approximation page:214 type:ordinary-paragraph
{In what follows, we consider using the Arnoldi process in Section~\ref{subsec:arnoldi-process} to compute basis vectors of $\mathcal{K}_m(A,\boldsymbol{x}_0)$ and give an approximate solution of \eqref{eq:matrix-exp-krylov-projection}. From the Arnoldi process in Section~\ref{subsec:arnoldi-process}, it follows from $\boldsymbol{x}_0=\|\boldsymbol{x}_0\|V_m\boldsymbol{e}_1$ that we obtain $\boldsymbol{x}_m=\|\boldsymbol{x}_0\|V_m(V_m^{\mathrm H}e^{tA}V_m)\boldsymbol{e}_1$, where $\boldsymbol{e}_1=[1,0,\ldots,0]^{\top}$. Furthermore, using the relation $V_m^{\mathrm H}AV_m=H_m$ (Hessenberg matrix) in \eqref{eq:arnoldi-projected-matrix}, we consider an approximation $V_m^{\mathrm H}e^{tA}V_m\approx e^{tV_m^{\mathrm H}AV_m}=e^{tH_m}$. Then we have}
\begin{equation}\label{eq:matrix-exp-arnoldi-approximation}\boldsymbol{x}_m=\|\boldsymbol{x}_0\|V_m(V_m^{\mathrm H}e^{tA}V_m)\boldsymbol{e}_1\approx\|\boldsymbol{x}_0\|V_me^{tH_m}\boldsymbol{e}_1.\end{equation}

% block: matrix-exponential-projected-cost page:214 type:ordinary-paragraph
{$e^{tH_m}\boldsymbol{e}_1$ is the multiplication of an $m\times m$ small matrix exponential function and a vector, whose computational costs are of $O(m^2)$. This comes from the fact that the multiplication of the Padé approximant (to $e^{tH_m}$) and a vector requires $H_m^{-1}\boldsymbol{v}$, i.e., solving linear systems $H_m\boldsymbol{z}=\boldsymbol{v}$, whose computational cost is of $O(m^2)$.}

% block: matrix-exp-hermitian-error-transition page:214 type:ordinary-paragraph
{If matrix $A$ is Hermitian, a suitable choice of $m$ for satisfying a given error tolerance is provided by using Theorem~\ref{thm:matrix-exp-krylov-error}.}

% block: matrix-exp-hermitian-error-theorem page:214 type:theorem
\Needspace{14\baselineskip}
\begin{Theorem}{\BookCite{99}, Theorem 2}\label{thm:matrix-exp-krylov-error}
{Let $A$ be Hermitian and all the eigenvalues lie in the interval $[-4\rho,0]$. Then the error of \eqref{eq:matrix-exp-arnoldi-approximation} is given by}
\[\|e^{tA}\boldsymbol{x}_0-\|\boldsymbol{x}_0\|_2V_me^{tH_m}\boldsymbol{e}_1\|\le\begin{cases}c\cdot e^{-\frac{m^2}{5\rho t}}&(\sqrt{4\rho t}\le m\le2\rho t),\\c\cdot\frac{1}{\rho te^{\rho t}}\cdot\left(\frac{e\rho t}{m}\right)^m&(2\rho t\le m),\end{cases}\]
{where $c=10\|\boldsymbol{x}_0\|$.}
\end{Theorem}

% block: matrix-exp-practical-dimension-selection page:214 type:ordinary-paragraph
{In practice, first estimate the minimum eigenvalue $\lambda_{\min}$ of $A$ and set $\rho=-\lambda_{\min}/4$. Then, the dimension of Krylov subspace $m$ is chosen so that the upper bound error in Theorem~\ref{thm:matrix-exp-krylov-error} is less than the given tolerance, e.g., $10^{-12}$. The usefulness of the error bound is shown in \BookCite{139}, together with devising a practical computation avoiding the loss of orthogonality regarding the basis vectors of Krylov subspaces and an efficient method of parallel computation.}
